% !TEX program = xelatex

\documentclass{resume}
\usepackage{xeCJK}
\setCJKmainfont[BoldFont=STSongti-SC-Black]{STSongti-SC-Regular}
\setCJKsansfont{STHeitiSC-Medium}
\setlist{itemsep=0.15em, topsep=0.3em, leftmargin=1.5pc}
\setlist[itemize,2]{label=\textbullet}
% 中英文标题统一使用正常字形，避免英文小型大写与中文字体回退
\titleformat{\section}{\normalfont\Large\raggedright}{}{0em}{}[\titlerule]
\renewcommand{\name}[1]{\centerline{\normalfont\Huge #1}\vspace{1.25ex}}
\clubpenalty=10000
\widowpenalty=10000
\raggedbottom

% 详细项目经历：呈现问题、实现与结果
\newcommand{\starblock}[4]{%
  \par\vspace{0.55em}%
  {\large\textbf{#1}}\par\vspace{0.12em}%
  \textbf{要解决的问题：}#2\par
  \textbf{做了什么：}#3\par
  \textbf{结果：}#4\par
}

\newcommand{\detailheading}[1]{\par\vspace{0.7em}{\large\textbf{#1}}\par\nopagebreak\vspace{0.2em}}
\newcommand{\pptpart}[1]{\par\vspace{0.7em}{\Large\textbf{#1}}\par\nopagebreak\vspace{0.2em}}
\begin{document}
\pagenumbering{arabic}
\pagestyle{plain}

\name{童浩}

\basicInfo{
  (+86) 188-5098-6208 \textperiodcentered\
  \href{mailto:Z1310554512@163.com}{Z1310554512@163.com} \textperiodcentered\
  中共党员
}
\centerline{\href{https://github.com/EmiyaredA}{\faGithub\ github.com/EmiyaredA}}

\section{个人摘要}
商汤科技算法研究员，工作覆盖 Agent 应用、知识检索与办公智能化。具备端到端演示系统设计、前后端集成和文档处理经验，围绕客服服务流程、PPT 生成与编辑恢复、笔记助手及训练数据构建开展开发；结合原文证据、规则校验和人工审核支持质量迭代。

\section{我会什么}
\begin{itemize}
  \item \textbf{业务：}智能客服、知识库问答、联网搜索、办公文档生成（表格 / 笔记 / PPT）。
  \item \textbf{基础：}Python、JavaScript / TypeScript、SQL；Dify 工作流编排、向量与 BM25 混合检索、Agent 工具调用、OCR 与文档重建。
  \item \textbf{用 AI 开发：}模型生成与规则校验结合，执行结果交叉对拍；保留原文证据与中间产物，以版本快照、差异审核和回归评测控制变更。
\end{itemize}

\section{项目一览}
\datedsubsection{\textbf{AI 客服中台}，独立开发}{2026.08 -- 至今}
\begin{itemize}
  \item 独立完成端到端演示系统的设计与搭建，实现配置发布、知识检索、工具与工作流、实时语音及人工交接，支持会话驱动优化与纯检索评测。
\end{itemize}

\datedsubsection{\textbf{商汤科技}，算法研究员}{2026.07 -- 至今}
\datedline{算法研究员（实习）}{2024.08 -- 2026.06}
\begin{itemize}
  \item \textbf{PPT：}完成项目演示系统的设计与搭建，实现结构化模板、HTML 模板、整页生图与重建三种方案，贯通内容规划、页面生成与可编辑 PPTX 交付。
  \item \textbf{笔记：}编排 Dify 查询改写与联网检索，实现资料筛选、编号引用、文档问答及选区改写与续写的上下文和输出适配；支持笔记格式转换与导入。
  \item \textbf{数据工程｜样本与标注：}补齐种子数据集中全部表格类型、分布均衡的训练样本，并用标注网站做质量把关。
  \item \textbf{数据工程｜评测与校验：}自动评测与人工评分一致性 92.4\%，日均节省评测人力超过 4 小时；推理结果验证准确率由 85.7\% 提升至 97.9\%。
  \item \textbf{数据工程｜生成质量：}字段命中率由 70.3\% 提升至 89.6\%，整体生成准确率提升 8.4\%。
\end{itemize}

\section{教育与奖项}
\datedsubsection{\textbf{厦门理工学院}，计算机技术，工学硕士}{2023.09 -- 2026.06}
\datedsubsection{\textbf{福州大学}，计算机科学与技术，工学学士}{2018.09 -- 2022.06}
排名 20/151（前 15\%）。福州大学数学与计算机科学学院学生会副主席、校学生会网络科技部部长；国家励志奖学金、一等奖学金、三好学生、优秀共青团员。

第五届全国大学生集成电路创新创业大赛决赛\textbf{全国二等奖}（2021）；第十一届中国大学生服务外包创新创业大赛\textbf{全国三等奖}（2020）。

\newpage
\section{项目经历｜AI 客服中台（配置与知识）}
\textbf{技术栈：}Python、FastAPI、PostgreSQL、Chroma、React、TypeScript、WebSocket。

\textbf{系统架构：}Studio、网站嵌入页与浏览器扩展共用运行时客户端；PostgreSQL 保存业务状态，Chroma 为可重建的检索投影，按正式交付要求维护数据一致性与配置版本。

\detailheading{机器人配置、草稿试用与不可变发布}
\begin{itemize}
  \item \textbf{配置建模：}组织角色、服务范围、语气与禁止事项，绑定知识库、HTTP 工具与工作流；以统一配置模型管理配置来源，依据 OpenAPI 接口定义生成管理端类型，减少多入口配置差异。
  \item \textbf{AI 辅助构建：}将自然语言需求转为可核对的配置修改，读取当前草稿与历史，只合并所需字段并同步表单；构建结果先进入草稿，经文本或语音试用后再发布。
  \item \textbf{发布校验：}草稿试用后检查模型、知识成员就绪状态、工具引用和工作流有效性；按选择的检索模式校验知识依赖。
  \item \textbf{运行快照：}同一事务提交版本、知识关联与运行数据，渠道只读取不可变发布快照。写入与发布时校验配置版本，冲突时返回当前快照，由使用者处理差异。
\end{itemize}

\detailheading{企业资料处理：正文、分块与检索表示}
\begin{itemize}
  \item \textbf{正文与版本：}保留原始文件、来源内容哈希、完整正文及抽取版本；读取时通过内容哈希校验正文与源文件是否一致，过期或缺失显式报错。解析、分块、摘要与持久化分别记录进度，通过任务版本校验防止过期任务覆盖最新结果。
  \item \textbf{多种分块：}支持固定滑窗、递归、标题结构、整篇及 Q\&A。标题树使用 DFS 收集路径，子块携带祖先标题，超长内容按 token 分窗并重附标题，减少上下文丢失。Q\&A 可从模板解析或由模型抽取，派生问答保留原段关联。
  \item \textbf{检索表示：}分离原始正文与检索用文本，支持正文、模型摘要和自定义检索文本；按问题关键词、检索文本、原文的优先级选择向量输入，返回时使用原文。按权重组合文件名与内容向量，查询与建库使用同一 Embedding 模型。
  \item \textbf{可重建索引：}Chroma 保存向量、原文与成员版本元数据，更新后使 BM25 缓存失效；外部向量操作与数据库事务分离，写入前后复核资料是否仍有效、任务版本是否仍匹配，清理失效投影。
\end{itemize}

\detailheading{混合检索与生成式检索}
\begin{itemize}
  \item \textbf{混合路径：}融合向量与 BM25 候选，支持 RRF 排序或归一分数加权；核对资料就绪状态与处理批次版本，按范围去重截断，返回原文片段与来源。
  \item \textbf{生成式路径：}基于完整正文、目录与行号索引运行，无需先分块或向量化。Agent 浏览目录、搜索原文并按需扩读；服务端登记已展示证据，提交时校验证据来源、内容哈希和行列边界，只接受已读原文内的有效范围，无相关证据可返回空结果。
  \item \textbf{探索控制：}限制探索轮数、搜索与单次阅读字符预算，支持游标续读；证据充分时提前提交，预算耗尽后只开放证据选择，对可修复参数错误提供有限纠错重提。
  \item \textbf{访问范围：}按知识库成员状态限定检索，外部资料还需确认当前有效可用；接入成功、处理完成、成员就绪与真实召回分别验证，避免把索引状态当作回答正确性的证明。
\end{itemize}

\newpage
\section{项目经历｜AI 客服中台（运行与迭代）}
\detailheading{客服 Agent 服务回合与事务协调}
\begin{itemize}
  \item \textbf{上下文装配：}读取已发布配置、近期对话、用户信息与附件，按需选择知识检索、HTTP 工具及工作流；工具结果返回模型，支持多轮调用后整合回答与引用。
  \item \textbf{工具与流程控制：}支持业务问题先查知识、闲聊直接回应；按渠道、权限和剩余次数开放工具，相同检索参数复用结果，失败或无结果也计次。工作流定义节点任务、工具、转人工与跳转条件，服务端校验转移，持久化进度并防止并发覆盖；会话固定发布版本，轮次超限按配置跳转、失败或转人工。
  \item \textbf{回合提交：}REST 与 SSE 共用回合协调器，先短事务保存用户消息，在事务外等待模型、网络和工具，再在结果检查点结算助手回复、证据、工具与工作流记录，避免长事务持有业务状态。
  \item \textbf{取消与副作用：}断线或取消只阻止尚未开始的检查点；已发生的外部业务动作保留真实记录，按工具自身幂等、租约与错误规则处理，不能依赖数据库回滚撤销。
\end{itemize}

\detailheading{实时语音与人工交接}
\begin{itemize}
  \item \textbf{语音分层：}浏览器负责采集、VAD、上行音频与播放；后端协调上下文、工具与生命周期，通过适配层处理 DashScope、SenseAudio 的协议与事件差异，按供应商能力开放取消、恢复、工具和工作流。
  \item \textbf{插话处理：}疑似近端人声先触发可恢复播放暂停；误判回声时从暂停位置继续，暂停后仍确认为人声才通知供应商打断，分别管理暂停、确认与恢复状态。
  \item \textbf{人工接待单：}管理创建、排队、认领、完成与取消。文本存在活动接待单时阻止 AI 回合，结束后恢复；语音在排队等待时继续 AI 服务，坐席接单后取消 AI 通话并清理会话资源，交接历史消息、工具结果和工作流状态。
  \item \textbf{运行与接入：}通过持久任务推进邀约、心跳与恢复；网站嵌入由接入方后端保管渠道密钥，页面换取短时访问令牌。语音测试资料覆盖 80 个语义脚本、139 条干净音频和 200 个分层场景，供专项验收使用。
\end{itemize}

\detailheading{基于会话证据的配置优化}
\begin{itemize}
  \item 从已选会话和具体消息整理证据与目标，生成人设、服务边界、FAQ 和协作话术等内容建议，展示修改前后值、理由和消息引用，逐项审核、编辑与采纳。
  \item 依据字段白名单、保留规则和配置版本校验，将建议应用到草稿；保护已有 FAQ、禁止话题与转人工条件，不修改模型、知识库和工具绑定，冲突时保留修改并返回差异。
  \item 应用后使用固定会话集进行草稿复测，再由人工发布；建议与上线分开，保留会话证据、修改理由和审核结果，便于追溯迭代。
\end{itemize}

\detailheading{服务质量评估设计与纯检索评测}
\begin{itemize}
  \item \textbf{质量维度：}从名词沿用错误、重复检索、重复念词与异常重音等问题出发，梳理理解与纠正、对话交互、回答与语音表达、任务执行与可靠性，结合安全、时延、成本和客户体验选择指标。
  \item \textbf{分层设计：}区分组件、语音链路、Agent 任务与产品效果，分别关注 WER / CER、MOS、可听响应时延、打断恢复、任务成功率及多次运行均成功率、首次联系解决率与满意度；以音频、转写、统一时间轴、工具记录和最终业务状态关联证据。该部分为评估设计，未作为已取得的效果成绩。
  \item \textbf{纯检索评测：}已实现独立于功能测试的检索评测，维护版本化问题集与参考原文证据；生成式和向量路径对每例使用相同三条检索词，通过统一接口获取结果，按检索结果的唯一标识去重。
  \item \textbf{评分统计：}统计证据召回、证据精度、全部证据命中及不可答问题返回空结果的准确性；汇总分位数与成对差值，用固定种子 Bootstrap 分析波动。
\end{itemize}

\newpage
\section{项目经历｜PPT Agent}
\textbf{项目背景：}将用户主题与参考资料转为可讲述、可修改的演示文稿，重点处理内容与版式匹配、生成过程追踪以及图像页面的编辑能力恢复。

\textbf{交付范围：}内容规划、模板处理、逐页生成、渲染检查及分层重建；按各入口集成模块，形成可运行的端到端演示流程。

\pptpart{一、结构化模板}
以已有 PPTX 为模板，解析版式与内容角色，再根据用户需求选择模板、填充结构化内容，交付保留原生对象和样式的 PPTX。

\detailheading{模板预处理：结构解析与语义识别}
\begin{itemize}
  \item \textbf{PPTX 结构解析：}提取元素标识、类型、位置、尺寸与文字样式，整理版式及嵌套关系，为后续内容填充保留原始元素映射。
  \item \textbf{多模态语义识别：}结合页面截图识别页面类型、标题与副标题、主视觉及内容角色，补充仅靠物理结构无法表达的元素语义和关系。
  \item \textbf{模板归一化：}将解析得到的结构与语义信息合并成统一记录，保存为模板 JSON / 数据库目录；规则覆盖目前已知的幻灯片类型，并支持用户上传的模板。
\end{itemize}

\detailheading{内容生成：逐页规划、模板匹配与导出}
\begin{itemize}
  \item \textbf{生成页面任务：}根据用户需求生成大纲和逐页结构化任务，再按页面类型筛选候选模板，减少生成内容与版式类型不匹配的情况。
  \item \textbf{选择与填充：}生成模板选择结果及结构化内容，把标题、要点等内容映射回原始元素。通过文字容量提示约束内容长度，保留原生对象与模板样式。
  \item \textbf{并发编排：}编排逐页并发生成任务，再按页面顺序归集结果，保证导出顺序与大纲一致。
  \item \textbf{交付结果：}支持套用用户模板；大纲与幻灯片结构生成控制在 2 分钟以内。长文本、换行和复杂元素仍通过渲染检查确认最终效果。
\end{itemize}

\pptpart{二、HTML 模板}
将视觉设计转为可复用的 HTML / CSS 模板，结合检索语料生成页面，通过浏览器渲染检查；交付原生 PPTX 时再转换，并核对浏览器呈现与导出结果。
\begin{itemize}
  \item \textbf{模板生产：}将单页视觉设计或网格拆分后的设计转为 HTML，提取页面结构定义与样式信息，经过浏览器预览后沉淀为模板，支持持续扩充版式。
  \item \textbf{内容准备：}围绕主题执行搜索与语料整理，组织大纲及页面结构，按结构匹配 HTML 模板。对内容页取出其指定文档知识块，作为外部语料与模板共同送入页面生成模块。
  \item \textbf{渲染与检查：}实现 HTML 页面生成、截图与结构提取；使用页面边界检查辅助发现溢出，并结合独立视觉审阅模块检查呈现。模块按具体入口组合，保留各阶段输入输出供定位问题。
\end{itemize}

\newpage
\section{项目经历｜PPT Agent（续）}
\pptpart{三、图像重建}
为恢复整页图像的文字与版式编辑能力，以生成的视觉图、已有图像或 PDF 为输入，识别布局并分离背景、文字和视觉素材，保留位置、样式与层级，合成为原生文本框及独立图片组成的 PPTX。

\detailheading{整页视觉生成与任务追踪}
\begin{itemize}
  \item \textbf{内容规划：}把用户主题、受众、页数、风格与重点组织为演示大纲和逐页内容，保存结构化大纲用于按页生成完整视觉图；该链路的大纲当前基于用户需求，尚未使用检索补充资料。
  \item \textbf{按页生成与归集：}应用统一风格组织页面任务，处理缓存、失败检查与重试，按页序归集图像并拼接 PDF，再进入逐页重建；已有 PDF 或图像也可直接进入重建流程。
  \item \textbf{任务记录：}使用统一任务标识关联用户请求、工作目录、生成图、重建结果和日志，保留各阶段输入输出与失败信息，便于重试、断点处理和效果对照。
\end{itemize}

\detailheading{布局识别与元素表}
接入 Paddle 布局识别，提取标题、正文、列表和图像等区域的类别、像素坐标与层级，形成元素表。依次执行区域裁剪、遮罩构建、背景生成、内容与样式提取和单页渲染；不同产物分路处理后再合成。

\detailheading{背景、文字与视觉素材的分层恢复}
\begin{itemize}
  \item \textbf{背景恢复：}根据文字与视觉区域生成移除遮罩，向外膨胀边界以覆盖残留边缘，再使用 LaMa 修复被遮住的背景纹理，作为最终页面底层。
  \item \textbf{文字恢复：}对裁剪区域逐元素 OCR，提取字体、字号、颜色等样式，使用最多 5 个线程提取样式，生成原生文本框；识别失败时回退为图片，保留可用页面。
  \item \textbf{视觉区域处理：}按检测框裁剪独立图片；对更高层元素造成的重叠，用背景对应区域替换，缺少背景时调用局部修补，减少文字等上层对象重复出现在底图中的问题。
  \item \textbf{坐标和层级合成：}将像素坐标映射为 PPTX 的 EMU 坐标，按元素层级组织前后关系，合成修复背景、原生文字和独立视觉图片，输出可继续手动改稿的 PPTX。
\end{itemize}

\section{PPT Agent 项目总结}
\begin{itemize}
  \item \textbf{结构化模板：}适合固定风格、重复汇报和频繁在 PPT 内改稿，复用原生对象与样式；页面表达受模板覆盖、内容容量和换行约束。
  \item \textbf{HTML 模板：}适合通过 HTML / CSS 维护和扩充代码版式，便于浏览器渲染检查；交付原生 PPTX 需要转换。
  \item \textbf{图像重建：}适合定制视觉与图文叙事，以及已有图像或 PDF 的编辑恢复；文字恢复为原生对象，复杂图表和插画通常保留为独立矩形图片。
\end{itemize}

\newpage
\section{项目经历｜笔记 AI}
\textbf{项目背景：}用户会结合当前文档或选区问答、改写、续写，也会连续追问外部资料；需要补全问题、组织文档上下文，并统一不同来源的资料与任务输出。

\textbf{个人工作：}编排 Dify 搜索工作流，实现 Python 节点、筛选排序和编号引用；组织文档、选区与指令，适配问答、改写、续写和搜索任务的上下文及输出；另参与笔记格式转换与导入。

\detailheading{多轮查询改写与搜索路由}
\begin{itemize}
  \item \textbf{补全追问背景：}将用户问题与历史对话组织为工作流输入，把“相关论文有哪些”等依赖前文的问题改写成可独立检索的查询，保留当前讨论主题。
  \item \textbf{识别与校验类型：}校验显式指定的网页、学术、科学、视频或社交类型，未指定时自动判断；学术调用 arXiv，其余四类使用 SearXNG，无需搜索时返回空资料。
  \item \textbf{组织节点依赖：}先获取并统一资料，再选择处理模式；相似度重排分支依次计算查询向量、文档向量，再执行相关性筛选与排序。
  \item \textbf{网页子流程设计：}设计网页初筛、正文抓取、分块摘要与向量过滤，作为网页资料处理的独立子流程。
\end{itemize}

\detailheading{候选资料处理与两种排序策略}
\begin{itemize}
  \item \textbf{统一资料输入：}组织摘要或搜索结果片段与已有来源元数据，过滤空内容；不将搜索片段视为已抓取全文或视频字幕。文档向量使用迭代节点，并行数为 10。
  \item \textbf{原序筛选：}保留来源顺序，按输出上限截取资料，不计算向量相似度，适合按原始召回顺序交付资料的调用方式。
  \item \textbf{相似度重排：}计算查询与文档的余弦相似度，保留严格大于阈值 0 的候选，降序排序后截取。两种策略为互斥分支。
  \item \textbf{来源上限：}arXiv 最多获取 5 篇论文，网络取前 20 条候选；两种策略最终最多返回 15 条，实际数量取决于来源与过滤，三个数量为不同阶段的独立约束。
\end{itemize}

\detailheading{直接回答、工具输出与上层助手适配}
\begin{itemize}
  \item \textbf{直接回答工作流：}将筛选后的正文按编号组织为上下文，结合用户问题与系统提示生成回答，同时返回原始资料，保留回答与来源的对应关系。
  \item \textbf{搜索工具工作流：}将资料正文、搜索类型与专用提示组织为结构化工具结果，交给上层调用方，支持由助手继续组织回答。
  \item \textbf{Python 适配层：}解析结构化资料，为每篇文档添加编号和来源标记，保留正文和已有元信息，生成格式化搜索结果、系统提示与用户提示，使答案可以引用编号来源。
\end{itemize}

\detailheading{文档助手：上下文组织与任务输出}
\begin{itemize}
  \item \textbf{上下文与分派：}组合当前文档、精确选区、用户指令、日期与任务提示，保留文档结构和选区位置，分派文档问答、选区改写、内容续写与联网搜索；按需将编号资料补入回答上下文。
  \item \textbf{输出适配：}问答输出回答内容，改写输出覆盖选区的替换内容，续写输出插入位置后的新增内容；约定标题、列表、表格、代码、公式和引用的结构化表达，并清理输出代码围栏。页面替换与插入由宿主执行。
\end{itemize}

\detailheading{笔记格式转换与迁移}
将模型输出映射为前端 HTML；解析多种主流笔记格式，统一为内部结构，在导入时保留正文与元信息，支持既有资料迁移到笔记应用。

\newpage
\section{项目经历｜数据工程}

\starblock{Excel 样本扩充与标注协作}{%
  训练数据需要覆盖不同表格类型与业务分布，标注结果需要可执行、可复查。%
}{%
  以种子集为基准，人工设计表头和数据行骨架，结合大语言模型与规则模板脚本泛化教育、金融等业务内容，通过 Python 导出 Excel，并覆盖排序、筛选和修改等任务。对原始 Prompt 继续泛化，形成 Excel、Prompt、JSA 对应样本。标注网站支持多角色协作，对同一内容多次采样并交叉对拍，自动筛出异常并进入人工复审。%
}{%
  训练样本覆盖种子集全部表格类型，保持业务分布均衡；建立生成、标注、异常筛查与复审流程。%
}

\starblock{JSA 执行对拍与自动评测}{%
  代码文本不同不代表执行结果不同；直接比较文件也容易受到冗余信息影响。人工评分耗时且口径不一。%
}{%
  交叉执行人工标注和具备 JSA 基础知识的自定义 GPT 输出，在相同输入上生成多份 Excel。将结果读取为 Python 对象，剔除冗余信息，序列化后计算 MD5 比较规范化内容，并结合语义比对定位差异。按结构化维度建立加权评分，用人工评分修正口径，用于 A/B 验证与回归。%
}{%
  自动评测与人工评分一致性达到 92.4\%，日均节省评测人力超过 4 小时；推理结果验证准确率由 85.7\% 提升至 97.9\%。%
}

\starblock{GRPO 输出校验与结构化奖励}{%
  输出容易缺失关键字段；文本、日期、数值的格式、精度与单位差异会干扰正确性判定。%
}{%
  建立文本、日期、数值的提取优先级：先完全匹配文本，再提取并验证日期值、统一格式，最后按符号、数值、百分比与单位解析数值表达式。按较低精度进行容差匹配；带单位时先完整匹配，失败后按业务规则去除单位重试；暂时停用与纯数值冲突的紧凑日期模板。在 GRPO 中另加入基于正则的分层关键词奖励，约束业务所需结构化字段。%
}{%
  字段命中率由 70.3\% 提升至 89.6\%，整体生成准确率提升 8.4\%；结构缺失与格式错误下降，错误案例总体占比减少 10\%，收敛更快。%
}

\end{document}
